\IfFileExists{stacks-project.cls}{%
\documentclass{stacks-project}
}{%
\documentclass{amsart}
}

% For dealing with references we use the comment environment
\usepackage{verbatim}
\newenvironment{reference}{\comment}{\endcomment}
%\newenvironment{reference}{}{}
\newenvironment{slogan}{\comment}{\endcomment}
\newenvironment{history}{\comment}{\endcomment}

% For commutative diagrams we use Xy-pic
\usepackage[all]{xy}

% We use 2cell for 2-commutative diagrams.
\xyoption{2cell}
\UseAllTwocells

% We use multicol for the list of chapters between chapters
\usepackage{multicol}

% This is generally recommended for better output
\usepackage{lmodern}
\usepackage[T1]{fontenc}

% For cross-file-references
\usepackage{xr-hyper}

% Package for hypertext links:
\usepackage{hyperref}

% For any local file, say "hello.tex" you want to link to please
% use \externaldocument[hello-]{hello}
\externaldocument[introduction-]{introduction}
\externaldocument[conventions-]{conventions}
\externaldocument[sets-]{sets}
\externaldocument[categories-]{categories}
\externaldocument[topology-]{topology}
\externaldocument[sheaves-]{sheaves}
\externaldocument[sites-]{sites}
\externaldocument[stacks-]{stacks}
\externaldocument[fields-]{fields}
\externaldocument[algebra-]{algebra}
\externaldocument[brauer-]{brauer}
\externaldocument[homology-]{homology}
\externaldocument[derived-]{derived}
\externaldocument[simplicial-]{simplicial}
\externaldocument[more-algebra-]{more-algebra}
\externaldocument[smoothing-]{smoothing}
\externaldocument[modules-]{modules}
\externaldocument[sites-modules-]{sites-modules}
\externaldocument[injectives-]{injectives}
\externaldocument[cohomology-]{cohomology}
\externaldocument[sites-cohomology-]{sites-cohomology}
\externaldocument[dga-]{dga}
\externaldocument[dpa-]{dpa}
\externaldocument[sdga-]{sdga}
\externaldocument[hypercovering-]{hypercovering}
\externaldocument[schemes-]{schemes}
\externaldocument[constructions-]{constructions}
\externaldocument[properties-]{properties}
\externaldocument[morphisms-]{morphisms}
\externaldocument[coherent-]{coherent}
\externaldocument[divisors-]{divisors}
\externaldocument[limits-]{limits}
\externaldocument[varieties-]{varieties}
\externaldocument[topologies-]{topologies}
\externaldocument[descent-]{descent}
\externaldocument[perfect-]{perfect}
\externaldocument[more-morphisms-]{more-morphisms}
\externaldocument[flat-]{flat}
\externaldocument[groupoids-]{groupoids}
\externaldocument[more-groupoids-]{more-groupoids}
\externaldocument[etale-]{etale}
\externaldocument[chow-]{chow}
\externaldocument[intersection-]{intersection}
\externaldocument[pic-]{pic}
\externaldocument[weil-]{weil}
\externaldocument[adequate-]{adequate}
\externaldocument[dualizing-]{dualizing}
\externaldocument[duality-]{duality}
\externaldocument[discriminant-]{discriminant}
\externaldocument[derham-]{derham}
\externaldocument[local-cohomology-]{local-cohomology}
\externaldocument[algebraization-]{algebraization}
\externaldocument[curves-]{curves}
\externaldocument[resolve-]{resolve}
\externaldocument[models-]{models}
\externaldocument[functors-]{functors}
\externaldocument[equiv-]{equiv}
\externaldocument[pione-]{pione}
\externaldocument[etale-cohomology-]{etale-cohomology}
\externaldocument[proetale-]{proetale}
\externaldocument[relative-cycles-]{relative-cycles}
\externaldocument[more-etale-]{more-etale}
\externaldocument[trace-]{trace}
\externaldocument[crystalline-]{crystalline}
\externaldocument[spaces-]{spaces}
\externaldocument[spaces-properties-]{spaces-properties}
\externaldocument[spaces-morphisms-]{spaces-morphisms}
\externaldocument[decent-spaces-]{decent-spaces}
\externaldocument[spaces-cohomology-]{spaces-cohomology}
\externaldocument[spaces-limits-]{spaces-limits}
\externaldocument[spaces-divisors-]{spaces-divisors}
\externaldocument[spaces-over-fields-]{spaces-over-fields}
\externaldocument[spaces-topologies-]{spaces-topologies}
\externaldocument[spaces-descent-]{spaces-descent}
\externaldocument[spaces-perfect-]{spaces-perfect}
\externaldocument[spaces-more-morphisms-]{spaces-more-morphisms}
\externaldocument[spaces-flat-]{spaces-flat}
\externaldocument[spaces-groupoids-]{spaces-groupoids}
\externaldocument[spaces-more-groupoids-]{spaces-more-groupoids}
\externaldocument[bootstrap-]{bootstrap}
\externaldocument[spaces-pushouts-]{spaces-pushouts}
\externaldocument[spaces-chow-]{spaces-chow}
\externaldocument[groupoids-quotients-]{groupoids-quotients}
\externaldocument[spaces-more-cohomology-]{spaces-more-cohomology}
\externaldocument[spaces-simplicial-]{spaces-simplicial}
\externaldocument[spaces-duality-]{spaces-duality}
\externaldocument[formal-spaces-]{formal-spaces}
\externaldocument[restricted-]{restricted}
\externaldocument[spaces-resolve-]{spaces-resolve}
\externaldocument[formal-defos-]{formal-defos}
\externaldocument[defos-]{defos}
\externaldocument[cotangent-]{cotangent}
\externaldocument[examples-defos-]{examples-defos}
\externaldocument[algebraic-]{algebraic}
\externaldocument[examples-stacks-]{examples-stacks}
\externaldocument[stacks-sheaves-]{stacks-sheaves}
\externaldocument[criteria-]{criteria}
\externaldocument[artin-]{artin}
\externaldocument[quot-]{quot}
\externaldocument[stacks-properties-]{stacks-properties}
\externaldocument[stacks-morphisms-]{stacks-morphisms}
\externaldocument[stacks-limits-]{stacks-limits}
\externaldocument[stacks-cohomology-]{stacks-cohomology}
\externaldocument[stacks-perfect-]{stacks-perfect}
\externaldocument[stacks-introduction-]{stacks-introduction}
\externaldocument[stacks-more-morphisms-]{stacks-more-morphisms}
\externaldocument[stacks-geometry-]{stacks-geometry}
\externaldocument[moduli-]{moduli}
\externaldocument[moduli-curves-]{moduli-curves}
\externaldocument[examples-]{examples}
\externaldocument[exercises-]{exercises}
\externaldocument[guide-]{guide}
\externaldocument[desirables-]{desirables}
\externaldocument[coding-]{coding}
\externaldocument[obsolete-]{obsolete}
\externaldocument[fdl-]{fdl}
\externaldocument[index-]{index}

% Theorem environments.
%
\theoremstyle{plain}
\newtheorem{theorem}[subsection]{Theorem}
\newtheorem{proposition}[subsection]{Proposition}
\newtheorem{lemma}[subsection]{Lemma}

\theoremstyle{definition}
\newtheorem{definition}[subsection]{Definition}
\newtheorem{example}[subsection]{Example}
\newtheorem{exercise}[subsection]{Exercise}
\newtheorem{situation}[subsection]{Situation}

\theoremstyle{remark}
\newtheorem{remark}[subsection]{Remark}
\newtheorem{remarks}[subsection]{Remarks}

\numberwithin{equation}{subsection}

% Macros
%
\def\lim{\mathop{\mathrm{lim}}\nolimits}
\def\colim{\mathop{\mathrm{colim}}\nolimits}
\def\Spec{\mathop{\mathrm{Spec}}}
\def\Hom{\mathop{\mathrm{Hom}}\nolimits}
\def\Ext{\mathop{\mathrm{Ext}}\nolimits}
\def\SheafHom{\mathop{\mathcal{H}\!\mathit{om}}\nolimits}
\def\SheafExt{\mathop{\mathcal{E}\!\mathit{xt}}\nolimits}
\def\Sch{\mathit{Sch}}
\def\Mor{\mathop{\mathrm{Mor}}\nolimits}
\def\Ob{\mathop{\mathrm{Ob}}\nolimits}
\def\Sh{\mathop{\mathit{Sh}}\nolimits}
\def\NL{\mathop{N\!L}\nolimits}
\def\CH{\mathop{\mathrm{CH}}\nolimits}
\def\proetale{{pro\text{-}\acute{e}tale}}
\def\etale{{\acute{e}tale}}
\def\QCoh{\mathit{QCoh}}
\def\Ker{\mathop{\mathrm{Ker}}}
\def\Im{\mathop{\mathrm{Im}}}
\def\Coker{\mathop{\mathrm{Coker}}}
\def\Coim{\mathop{\mathrm{Coim}}}

% Boxtimes
%
\DeclareMathSymbol{\boxtimes}{\mathbin}{AMSa}{"02}

%
% Macros for moduli stacks/spaces
%
\def\QCohstack{\mathcal{QC}\!\mathit{oh}}
\def\Cohstack{\mathcal{C}\!\mathit{oh}}
\def\Spacesstack{\mathcal{S}\!\mathit{paces}}
\def\Quotfunctor{\mathrm{Quot}}
\def\Hilbfunctor{\mathrm{Hilb}}
\def\Curvesstack{\mathcal{C}\!\mathit{urves}}
\def\Polarizedstack{\mathcal{P}\!\mathit{olarized}}
\def\Complexesstack{\mathcal{C}\!\mathit{omplexes}}
% \Pic is the operator that assigns to X its picard group, usage \Pic(X)
% \Picardstack_{X/B} denotes the Picard stack of X over B
% \Picardfunctor_{X/B} denotes the Picard functor of X over B
\def\Pic{\mathop{\mathrm{Pic}}\nolimits}
\def\Picardstack{\mathcal{P}\!\mathit{ic}}
\def\Picardfunctor{\mathrm{Pic}}
\def\Deformationcategory{\mathcal{D}\!\mathit{ef}}
\usepackage{cleveref}
\crefname{originalreference}{original reference}{original references}
\Crefname{originalreference}{Original reference}{Original references}
\makeatletter
\newlabel{algebra-lemma-cokernel-flat}{{00O0}{source}{}{sourcepointer-1}{}}
\newlabel{algebra-lemma-cokernel-flat@cref}{{[originalreference][0][]00O0}{[1][0][]source}{}{}{}}
\makeatother
\crefname{originallemma}{original lemma}{original lemmas}
\Crefname{originallemma}{Original lemma}{Original lemmas}
\makeatletter
\newlabel{lemma-represent-FAd}{{05YP}{source}{}{sourcepointer-2}{}}
\newlabel{lemma-represent-FAd@cref}{{[originallemma][0][]05YP}{[1][0][]source}{}{}{}}
\makeatother
\crefname{originalreference}{original reference}{original references}
\Crefname{originalreference}{Original reference}{Original references}
\makeatletter
\newlabel{spaces-groupoids-definition-quotient-stack}{{044Q}{source}{}{sourcepointer-3}{}}
\newlabel{spaces-groupoids-definition-quotient-stack@cref}{{[originalreference][0][]044Q}{[1][0][]source}{}{}{}}
\makeatother
\crefname{originalreference}{original reference}{original references}
\Crefname{originalreference}{Original reference}{Original references}
\makeatletter
\newlabel{stacks-lemma-stackify-groupoids}{{02ZP}{source}{}{sourcepointer-4}{}}
\newlabel{stacks-lemma-stackify-groupoids@cref}{{[originalreference][0][]02ZP}{[1][0][]source}{}{}{}}
\makeatother
\title{The finite Hilbert stack of a point is the quotient stack of the free-algebra functor}
\author{The Stacks Project Authors}
\date{stacks-project:criteria.tex:local-snapshot-20261012; GNU FDL 1.2 or later}
\begin{document}
\maketitle
\noindent\textit{Editorial note.} The sentence 'We omit the verification that this construction is compatible with base change so we get indeed a functor' is retained verbatim; the omitted check is the routine functoriality of relative Spec under base change, flagged for the reviewer.\par\medskip
\noindent\textit{Allowed standard prerequisites (editorial).}\begin{itemize}
\item Quotient stack as stackification. spaces-groupoids-definition-quotient-stack (cross-chapter, named). Identifies [X/G] with the stackification of the action presheaf in groupoids.
\item Stackification criterion. stacks-lemma-stackify-groupoids (cross-chapter, named). Reduces the equivalence to the two displayed conditions.
\item Cokernel flatness. algebra-lemma-cokernel-flat (cross-chapter, named). Justifies using 1 as a local basis element.
\end{itemize}\medskip
\ifdefined\hypertarget\hypertarget{sourcepointer-1}{}\fi
\noindent\textit{Source reference.} Original reference 00O0: focused-seeds/geometry/S9-crit-finite-hilbert-point/ref-tags.txt, source lines 5--5. Navigation pointer only: the counted proof cites this statement by \textbackslash{}ref as a standing input; the deduction the counted proof has to make is carried out inside the counted span, and no part of that deduction is delegated to this pointer.\par\medskip
\ifdefined\hypertarget\hypertarget{sourcepointer-2}{}\fi
\noindent\textit{Source reference.} Original lemma 05YP: sources/stacks-criteria/pinned-original-source.tex, source lines 2447--2447. Navigation pointer only: the counted proof cites this statement by \textbackslash{}ref as a standing input; the deduction the counted proof has to make is carried out inside the counted span, and no part of that deduction is delegated to this pointer.\par\medskip
\ifdefined\hypertarget\hypertarget{sourcepointer-3}{}\fi
\noindent\textit{Source reference.} Original reference 044Q: focused-seeds/geometry/S9-crit-finite-hilbert-point/ref-tags.txt, source lines 6--6. Navigation pointer only: the counted proof cites this statement by \textbackslash{}ref as a standing input; the deduction the counted proof has to make is carried out inside the counted span, and no part of that deduction is delegated to this pointer.\par\medskip
\ifdefined\hypertarget\hypertarget{sourcepointer-4}{}\fi
\noindent\textit{Source reference.} Original reference 02ZP: focused-seeds/geometry/S9-crit-finite-hilbert-point/ref-tags.txt, source lines 7--7. Navigation pointer only: the counted proof cites this statement by \textbackslash{}ref as a standing input; the deduction the counted proof has to make is carried out inside the counted span, and no part of that deduction is delegated to this pointer.\par\medskip
\begin{proposition}
\label{proposition-finite-hilbert-point}
The stack $\mathcal{H}_d$ is equivalent to the quotient stack
$[X/G]$ described above. In particular $\mathcal{H}_d$ is an
algebraic stack.
\end{proposition}

\begin{proof}
Note that by
Groupoids in Spaces, Definition
\ref{spaces-groupoids-definition-quotient-stack}
the quotient stack $[X/G]$ is the stackification of the
category fibred in groupoids associated to the ``presheaf in groupoids''
which associates to a scheme $T$ the groupoid
$$
(X(T), G(T) \times X(T), s, t, c).
$$
Since this ``presheaf in groupoids'' is isomorphic to $FA_d$ by
Lemma \ref{lemma-represent-FAd}
it suffices to prove that the $\mathcal{H}_d$ is the stackification
of (the category fibred in groupoids associated to the
``presheaf in groupoids'') $FA_d$. To do this we first define a
functor
$$
\Spec : FA_d \longrightarrow \mathcal{H}_d
$$
Recall that the fibre category of $\mathcal{H}_d$ over a scheme $T$
is the category of finite locally free morphisms $Z \to T$ of degree $d$.
Thus given a scheme $T$ and a free $d$-dimensional
$\Gamma(T, \mathcal{O}_T)$-algebra $m$ we may assign to this the object
$$
Z = \underline{\Spec}_T(\mathcal{A})
$$
of $\mathcal{H}_{d, T}$
where $\mathcal{A} = \mathcal{O}_T^{\oplus d}$ endowed with a
$\mathcal{O}_T$-algebra structure via $m$. Moreover, if $m'$ is
a second such free $d$-dimensional $\Gamma(T, \mathcal{O}_T)$-algebra
and if $\varphi : m \to m'$ is an isomorphism of these, then
the induced $\mathcal{O}_T$-linear map
$\varphi : \mathcal{O}_T^{\oplus d} \to \mathcal{O}_T^{\oplus d}$
induces an isomorphism
$$
\varphi : \mathcal{A}' \longrightarrow \mathcal{A}
$$
of quasi-coherent $\mathcal{O}_T$-algebras. Hence
$$
\underline{\Spec}_T(\varphi) :
\underline{\Spec}_T(\mathcal{A})
\longrightarrow
\underline{\Spec}_T(\mathcal{A}')
$$
is a morphism in the fibre category $\mathcal{H}_{d, T}$. We omit the
verification that this construction is compatible with base change so
we get indeed a functor $\Spec : FA_d \to \mathcal{H}_d$
as claimed above.

\medskip\noindent
To show that $\Spec : FA_d \to \mathcal{H}_d$ induces an equivalence
between the stackification of $FA_d$ and $\mathcal{H}_d$ it suffices to
check that
\begin{enumerate}
\item $\mathit{Isom}(m, m') = \mathit{Isom}(\Spec(m), \Spec(m'))$
for any $m, m' \in FA_d(T)$.
\item for any scheme $T$ and any object $Z \to T$ of $\mathcal{H}_{d, T}$
there exists a covering $\{T_i \to T\}$ such that $Z|_{T_i}$ is
isomorphic to $\Spec(m)$ for some $m \in FA_d(T_i)$, and
\end{enumerate}
see
Stacks, Lemma \ref{stacks-lemma-stackify-groupoids}.
The first statement follows from the observation that any isomorphism
$$
\underline{\Spec}_T(\mathcal{A})
\longrightarrow
\underline{\Spec}_T(\mathcal{A}')
$$
is necessarily given by a global invertible matrix $g$ when
$\mathcal{A} = \mathcal{A}' = \mathcal{O}_T^{\oplus d}$ as modules.
To prove the second statement let $\pi : Z \to T$ be a finite
locally free morphism of degree $d$. Then $\mathcal{A}$ is a locally
free sheaf $\mathcal{O}_T$-modules of rank $d$.
Consider the element $1 \in \Gamma(T, \mathcal{A})$. This element is
nonzero in $\mathcal{A} \otimes_{\mathcal{O}_{T, t}} \kappa(t)$
for every $t \in T$ since the scheme
$Z_t = \Spec(\mathcal{A} \otimes_{\mathcal{O}_{T, t}} \kappa(t))$
is nonempty being of degree $d > 0$ over $\kappa(t)$. Thus
$1 : \mathcal{O}_T \to \mathcal{A}$ can locally be used as the first
basis element (for example you can use
Algebra, Lemma \ref{algebra-lemma-cokernel-flat} parts (1) and (2)
to see this). Thus, after localizing on
$T$ we may assume that there exists an isomorphism
$\varphi : \mathcal{A} \to \mathcal{O}_T^{\oplus d}$
such that $1 \in \Gamma(\mathcal{A})$ corresponds to the first basis element.
In this situation the multiplication map
$\mathcal{A} \otimes_{\mathcal{O}_T} \mathcal{A} \to \mathcal{A}$
translates via $\varphi$ into a free $d$-dimensional algebra $m$ over
$\Gamma(T, \mathcal{O}_T)$. This finishes the proof.
\end{proof}

\end{document}
